%!TEX TS-program = xelatex
%!TEX encoding = UTF-8 Unicode
% Awesome CV LaTeX Template for CV / Cover Letter-style Research CV

%-------------------------------------------------------------------------------
% CONFIGURATIONS
%-------------------------------------------------------------------------------
\documentclass[11pt, a4paper]{awesome-cv}

\geometry{left=1.4cm, top=.8cm, right=1.4cm, bottom=1.8cm, footskip=.5cm}
\fontdir[fonts/]

\colorlet{awesome}{awesome-red}
\setbool{acvSectionColorHighlight}{true}
\renewcommand{\acvHeaderSocialSep}{\quad\textbar\quad}

%-------------------------------------------------------------------------------
% PERSONAL INFORMATION
%-------------------------------------------------------------------------------
\photo[circle,noedge,left]{profile}
\name{Changmin}{Lee}
\mobile{(+82) 10-9288-8981}
\email{lchm1106@unist.ac.kr}
\github{CkdalsKong}
\linkedin{changmin0lee}
\googlescholar{HV9YOPEAAAAJ}{Changmin Lee}
\homepage{https://changmin-lee.com}
\extrainfo{UAI Lab, UNIST}

%-------------------------------------------------------------------------------
\begin{document}

\makecvheader[R]

\makecvfooter
  {\today}
  {Changmin Lee~~~\cdotp~~~Curriculum Vitae}
  {}

%-------------------------------------------------------------------------------
% CV CONTENT
%-------------------------------------------------------------------------------
\begin{cvletter}

%-------------------------------------------------------------------------------
\lettersection{Education}

\textbf{Ulsan National Institute of Science and Technology (UNIST)} \\
M.S.-Ph.D. Integrated Program, Department of Computer Science and Engineering \\
UAI Lab, Advisor: Prof. Taesik Gong \hfill \textit{Mar. 2025 - Present}
\begin{itemize}[left=10pt]
    \item Research focus: Personalized AI, on-device RAG, efficient LLMs, and AI agents.
\end{itemize}

\textbf{Kyungpook National University} \\
B.S., Department of Computer Science and Engineering \hfill \textit{Mar. 2019 - Feb. 2025}
\begin{itemize}[left=10pt]
    \item Total GPA: \textbf{3.9/4.3}, Major GPA: \textbf{4.2/4.3}
\end{itemize}

%-------------------------------------------------------------------------------
\lettersection{Research Interests}
\begin{itemize}[left=10pt]
    \item Personalized AI, Retrieval-Augmented Generation, and long-term memory systems
    \item On-device AI, efficient LLMs, and resource-constrained inference
    \item AI agents, GUI agents, and memory-aware agent systems
    \item Preference-aware retrieval, memory construction, and adaptive AI systems
\end{itemize}

%-------------------------------------------------------------------------------
\lettersection{Research Experience}

\textbf{UAI Lab, UNIST} \\
Graduate Researcher \hfill \textit{Mar. 2025 - Present}
\begin{itemize}[left=10pt]
    \item Conduct research on personalized AI systems, on-device RAG, efficient memory construction, and AI agents.
    \item Developed preference-aligned memory construction methods for compact and efficient personalized retrieval on resource-constrained devices.
    \item Conducted experiments on retrieval efficiency, memory usage, latency, and generation quality across multiple personalized RAG benchmarks.
\end{itemize}

\textbf{Pattern Recognition \& Machine Intelligence Lab, Kyungpook National University} \\
Undergraduate Research Assistant \hfill \textit{Dec. 2023 - Sep. 2024}
\begin{itemize}[left=10pt]
    \item Conducted research on explainable AI, computer vision, medical image analysis, and machine learning.
    \item Developed a hybrid Transformer-CNN model for chest X-ray classification.
    \item Participated in regular research paper seminars and presented recent papers in computer vision and medical AI.
\end{itemize}

%-------------------------------------------------------------------------------
\lettersection{Publications and Preprints}
\begin{itemize}[left=10pt]
    \item \underline{Changmin Lee}, Jaemin Kim, Taesik Gong, ``From Volume to Value: Preference-Aligned Memory Construction for On-Device RAG,'' \textit{International Conference on Machine Learning (ICML)}, 2026. \textbf{Accepted}.
    \item \underline{Changmin Lee}, Taehwan Park, Jongwon Lee, Taesik Gong, ``$\Delta$Agent: Efficient Mobile GUI Agents via Transition-Guided Flexible Action Reuse,'' \textit{Preprint}, 2026. \textit{Under preparation}.
    \item \underline{Chang-min Lee}, Ho-kyung Shin, Woo-Jeoung Nam, ``Enhancing Chest X-Ray Classification with Multi-Class Token Hybrid Transformers,'' \textit{Journal of The Institute of Electronics and Information Engineers}, 2024.
    \item Dong-hyuk Kim, \underline{Chang-min Lee}, \textit{et al.}, ``Comparison of Recycling Waste Awareness Performance by Transfer Learning Model,'' presented at the \textit{Korean Multimedia Society}, 2024, Poster Session.
\end{itemize}

%-------------------------------------------------------------------------------
\lettersection{Awards and Fellowships}
\begin{itemize}[left=10pt]
    \item \textbf{NRF Master's Fellowship}, National Research Foundation of Korea, 2025.08 - 2026.09. \\
    Awarded the NRF Master's Fellowship, also known as \textit{Graduate Student Research Encouragement Grant for Master's Program}.
    \item \textbf{Excellence Award}, Undergraduate Research Support Program, \textit{Convergent Medical Scientist Training Program}, 2024.
\end{itemize}

%-------------------------------------------------------------------------------
\lettersection{Leadership and Service}

\textbf{UAI Lab, UNIST} \\
Lab Manager \hfill \textit{Jan. 2025 - Dec. 2025}
\begin{itemize}[left=10pt]
    \item Coordinated lab operations, internal communication, meeting schedules, and administrative tasks.
    \item Supported onboarding and communication among lab members for smooth research collaboration.
\end{itemize}

%-------------------------------------------------------------------------------
\lettersection{Teaching Experience}

\textbf{Teaching Assistant, Computer Networks (CSE351)} \\
Department of Computer Science and Engineering, UNIST \hfill \textit{Fall 2025}
\begin{itemize}[left=10pt]
    \item Supported course operations, announcements, grading, and student communication.
    \item Assisted students with course materials, assignments, and project-related questions.
\end{itemize}

\textbf{Teaching Assistant, On-Device AI for Smart Manufacturing (NOVA508)} \\
Department of Computer Science and Engineering, UNIST
\begin{itemize}[left=10pt]
    \item Supported project-based activities and on-device AI deployment practice.
    \item Assisted students with course materials, assignments, and hands-on implementation.
\end{itemize}

\textbf{Teaching Assistant, Data Preprocessing and Visualization} \\
AI Novatus Academia AI Practitioner Program (Gyeongnam), 6th Cohort \hfill \textit{May 2025}
\begin{itemize}[left=10pt]
    \item Supported theory lectures and hands-on exercises on data preprocessing and visualization.
\end{itemize}

\textbf{Teaching Assistant, Data Preprocessing and Visualization} \\
AI Novatus Academia AI Practitioner Program (Ulsan), 8th Cohort \hfill \textit{Jul. 2025}
\begin{itemize}[left=10pt]
    \item Supported theory lectures and hands-on exercises on data preprocessing and visualization.
\end{itemize}

\textbf{Teaching Assistant, Artificial Intelligence Fundamentals 2} \\
AI Novatus Academia AI Practitioner Program (Gyeongnam), 7th Cohort \hfill \textit{May - Jun. 2025}
\begin{itemize}[left=10pt]
    \item Supported theory lectures and student exercises for an introductory AI course.
\end{itemize}

\textbf{Teaching Assistant, Data Preprocessing} \\
AI Novatus Academia AI Practitioner Program (Gyeongnam), 7th Cohort \hfill \textit{May - Jun. 2025}
\begin{itemize}[left=10pt]
    \item Supported data preprocessing lectures and hands-on exercises.
\end{itemize}

\textbf{Teaching Assistant, Data Preprocessing and Visualization} \\
LG LDC Program \hfill \textit{Jul. 2025}
\begin{itemize}[left=10pt]
    \item Supported theory lectures and hands-on exercises on data preprocessing and visualization.
\end{itemize}

\textbf{Teaching Assistant, Data Preprocessing and Visualization} \\
LG LDC Program \hfill \textit{Jan. 2026}
\begin{itemize}[left=10pt]
    \item Supported theory lectures and hands-on exercises on data preprocessing and visualization.
\end{itemize}

%-------------------------------------------------------------------------------
\lettersection{Selected Research Projects}

\textbf{Preference-Aligned Memory Construction for On-Device RAG (EPIC)} \\
Graduate Researcher \hfill \textit{2025 - 2026}
\begin{itemize}[left=10pt]
    \item Designed an on-device memory construction framework that converts large volumes of raw user data into compact, preference-aligned memory.
    \item Evaluated the system across personalized RAG benchmarks with respect to accuracy, memory footprint, indexing latency, and retrieval latency.
    \item Accepted to \textit{ICML 2026}.
\end{itemize}

\textbf{$\Delta$Agent: Efficient Mobile GUI Agents via Transition-Guided Flexible Action Reuse} \\
Graduate Researcher \hfill \textit{2025 - Present}
\begin{itemize}[left=10pt]
    \item Developed transition-aware compact memory construction for mobile GUI agents with long-horizon task support.
    \item Investigated efficient action reuse guided by transition-level memory for resource-constrained mobile environments.
\end{itemize}

\textbf{Chest X-Ray Classification with Hybrid Transformer-CNN Model (MCTCheXFormer)} \\
Undergraduate Research Assistant \hfill \textit{Dec. 2023 - Oct. 2024}
\begin{itemize}[left=10pt]
    \item Developed a hybrid model combining CNN and Transformer-based modules for multi-label chest X-ray classification.
    \item Improved multi-label disease classification performance using multi-class tokens and token refinement mechanisms.
    \item Published in the \textit{IEIE Journal} and received an Excellence Award through the undergraduate research program.
\end{itemize}

\textbf{Comparison of Recycling Waste Awareness Performance} \\
Team Member \hfill \textit{Sep. 2023 - Dec. 2023}
\begin{itemize}[left=10pt]
    \item Presented a poster at the \textit{Korean Multimedia Society}, 2024.
    \item Compared recycling waste recognition models using transfer learning, including YOLOv8, YOLOv9, and RT-DETR.
    \item Analyzed the effectiveness of model architectures for automated recycling waste recognition.
\end{itemize}

%-------------------------------------------------------------------------------
\lettersection{AI Novatus and Industry-Academic Projects}

\textbf{Cutting Tool Dynamic Characteristics Prediction} \\
AI Novatus Academia (Gyeongnam), 6th Cohort PBL \hfill \textit{Jun. - Sep. 2025}
\begin{itemize}[left=10pt]
    \item Built a model to predict dynamic characteristics of machining tools from cutting-process images.
    \item Applied computer vision methods to manufacturing process analysis under the AI Novatus project-based learning program.
\end{itemize}

\textbf{Strawberry Smart Farm Growth Cycle Prediction} \\
AI Novatus Academia (Ulsan), 8th Cohort PBL \hfill \textit{Sep. - Dec. 2025}
\begin{itemize}[left=10pt]
    \item Developed a growth cycle classification model for smart-farm strawberry cultivation management.
    \item Applied classification methods to agricultural monitoring in the AI Novatus project-based learning program.
\end{itemize}

\textbf{Vacuum Cleaner Personalization Pipeline} \\
LG Electronics DIC Program PBL \hfill \textit{Aug. - Sep. 2025}
\begin{itemize}[left=10pt]
    \item Designed and implemented a personalization pipeline for LG vacuum cleaner user scenarios.
    \item Built an end-to-end pipeline for user-specific adaptation in consumer robotics applications.
\end{itemize}

\textbf{LG Electronics HS Future Technology Project} \\
Graduate Researcher \hfill \textit{Jun. - Sep. 2025}
\begin{itemize}[left=10pt]
    \item Developed a personalized model algorithm using Mixture-of-Experts (MoE) for user-specific adaptation.
    \item Conducted industry-academic collaboration on efficient personalization under resource constraints.
\end{itemize}

%-------------------------------------------------------------------------------
\lettersection{Other Project Experience}

\textbf{Recycling Waste Sorting Based on Transfer Learning} \\
Team Member \hfill \textit{Mar. 2024 - Jun. 2024}
\begin{itemize}[left=10pt]
    \item Participated in an industry-university linkage software project.
    \item Labeled a dataset of approximately 30 recycling waste categories using videos captured at a waste treatment plant.
    \item Fine-tuned YOLOv9, YOLOv8, and RT-DETR models through transfer learning and hyperparameter optimization.
    \item Proposed the best-performing YOLOv9 model to the company after comparing model performance.
\end{itemize}

\textbf{Smart Table Clock} \\
Team Member \hfill \textit{Sep. 2023 - Dec. 2023}
\begin{itemize}[left=10pt]
    \item Developed a smart table clock using an Arduino board to display widgets on an LCD screen.
    \item Implemented real-time bus arrival information using the Daegu bus system API.
    \item Built a web page to input data into the clock via Wi-Fi and integrated Arduino with network communication.
\end{itemize}

\textbf{Remote File Explorer} \\
Team Leader \hfill \textit{Mar. 2023 - Jun. 2023}
\begin{itemize}[left=10pt]
    \item Developed a remote file transfer system using C and implemented basic Linux-style file commands.
    \item Built a terminal-based remote file explorer using the curses library.
    \item Enabled file browsing, transfer, and management between macOS and Linux systems.
\end{itemize}

%-------------------------------------------------------------------------------
\lettersection{Skills}
\textbf{Programming} \quad Python, C, Java \\
\textbf{Deep Learning} \quad PyTorch, TensorFlow, Keras, scikit-learn \\
\textbf{RAG and Retrieval} \quad FAISS, dense retrieval, LLM-based evaluation, personalized RAG \\
\textbf{Systems and Tools} \quad Linux, Git, Docker, LaTeX, vLLM

\end{cvletter}

\end{document}
